\subsection{元素谱的分划}

设 A 是一个幺 Banach 代数，\(x \in A\)，且 \(K = \operatorname{Sp}_{\mathrm{A}}(x)\)。用 \(\Pi\) 表示 K 的在 K 中既开又闭的子集组成的集合。设 B 是由 x 生成的 A 的闭全子代数；它是交换的。

对每个 \(H \in \Pi\)，存在唯一的元素 \(f_{H}\)（属于 \(\mathcal{O}(K)\)），它在 H 的某邻域中等于 1，在 K---H 的某邻域中等于 0。我们令 \(j_{\mathrm{H}} = f_{\mathrm{H}}(x)\)。元素 \(j_{H}\) 是 A 的幂等元，称为与 x 和 H 相伴的幂等元，并且我们有下列公式：

\[
j _ {\mathrm{H} \cap \mathrm{H} ^ {\prime}} = j _ {\mathrm{H}} j _ {\mathrm{H} ^ {\prime}} = j _ {\mathrm{H} ^ {\prime}} j _ {\mathrm{H}}\tag{15}
\]

(H, H' ∈ Π)

\[
j _ {\mathrm {H\cup H^ {\prime}}} = j _ {\mathrm{H}} + j _ {\mathrm {H^ {\prime}}} - j _ {\mathrm {H^ {\prime}}} j _ {\mathrm{H}}\tag{16}
\]

\[
(\mathrm{H}, \mathrm{H} ^ {\prime} \in \Pi)
\]

\[
j _ {\emptyset} = 0, \qquad j _ {\mathrm{K}} = 1.
\]

设 \(H \in \Pi\)。我们定义 \(A_{H} = j_{H}Aj_{H}\)。这是 A 的一个闭子代数，以 \(j_{H}\) 为单位元（参见 I, p.~2 的引理 1）。我们还令 \(B_{H} = j_{H}Bj_{H}\) 以及 \(x_{H} = xj_{H} = j_{H}x = j_{H}xj_{H} \in B_{H}\)。

设 \(g_{H}\) 是 \(\mathcal{O}(K)\) 中由如下方式定义的元素：在 H 的邻域中 \(g_{\mathrm{H}}(z)=z\)，在 K---H 的邻域中 \(g_{\mathrm{H}}(z)=0\)。在 K 上有 \(g_{\mathrm{H}}(z)=f_{\mathrm{H}}(z)z\)，因此 \(x_{\mathrm{H}}=g_{\mathrm{H}}(x)\)。由此可得，若 \(H\neq K\)，则

\[
\mathrm{Sp} _ {\mathrm{A}} (x _ {\mathrm{H}}) = g _ {\mathrm{H}} (\mathrm{K}) = \mathrm{H} \cup \{0 \}.
\]

设 \(\lambda \in \mathbf{C} - \mathrm{H}\)。设 \(h_{\mathrm{H},\lambda}\) 是 \(\mathcal{O}(\mathrm{K})\) 中在 H 的某邻域中等于 \((\lambda - z)^{-1}\)、在 K-H 的某邻域中等于 0 的元素。我们有 \(h_{\mathrm{H},\lambda} = f_{\mathrm{H}}h_{\mathrm{H},\lambda}\) 以及 \((\lambda f_{\mathrm{H}} - g_{\mathrm{H}})h_{\mathrm{H},\lambda} = f_{\mathrm{H}}\)。若记 \(\mathrm{R_H}(x,\lambda) = h_{\mathrm{H},\lambda}(x)\)，则我们有 \(\mathrm{R_H}(x,\lambda) \in \mathrm{B_H}\)，并且

\[
\begin{array}{r} \mathrm {R_ {H}} (x, \lambda) (\lambda j _ {\mathrm{H}} - x _ {\mathrm{H}}) = (\lambda j _ {\mathrm{H}} - x _ {\mathrm{H}}) \mathrm {R_ {H}} (x, \lambda) = j _ {\mathrm{H}}, \\ \mathrm {R_ {H}} (x, \lambda) j _ {\mathrm {K-H}} = j _ {\mathrm {K-H}} \mathrm {R_ {H}} (x, \lambda) = 0. \end{array}\tag{17}
\]

特别地，\(\lambda \in \mathbf{C} - \mathrm{Sp}_{\mathrm{A_H}}(x_\mathrm{H})\)。

现在设 \(\lambda\in\mathrm{H}\)。假设 \(\lambda j_{\mathrm{H}}-x_{\mathrm{H}}\) 有逆 \(y\)（在 \(\mathrm{A}_{\mathrm{H}}\) 中）。利用公式 \(j_{\mathrm{H}}y=y\)（因为 \(y\in\mathrm{A}_{\mathrm{H}}\)）与 \(j_{\mathrm{K-H}}\mathrm{R}_{\mathrm{K-H}}(x,\lambda)=\mathrm{R}_{\mathrm{K-H}}(x,\lambda)\)（因为 \(\mathrm{R}_{\mathrm{K-H}}(x,\lambda)\in\mathrm{A}_{\mathrm{K-H}}\)），我们得到

\[
(\lambda - x) (y + \mathrm{R} _ {\mathrm{K-H}} (x, \lambda)) = (\lambda - x) (j _ {\mathrm{H}} y + j _ {\mathrm{K-H}} \mathrm{R} _ {\mathrm{K-H}} (x, \lambda)) =
\]

\[
(\lambda j _ {\mathrm{H}} - x j _ {\mathrm{H}}) y + (\lambda j _ {\mathrm {K-H}} - x j _ {\mathrm {K-H}}) \mathrm{R} _ {\mathrm {K-H}} (x, \lambda) = j _ {\mathrm{H}} + j _ {\mathrm {K-H}} = 1,
\]

（得益于应用于 K---H 的公式 (17)）。同样可以验证

\[
(y + \mathrm{R} _ {\mathrm{K-H}} (x, \lambda)) (\lambda - x) = 1.
\]

这证明 \(\lambda - x\) 在 A 中有逆 \(y + \mathrm{R}_{\mathrm{K - H}}(x, \lambda)\)，这是荒谬的。于是 \(\lambda \in \mathrm{Sp}_{\mathrm{A_H}}(x_{\mathrm{H}})\)。因此我们得出

\[
\mathrm{Sp} _ {\mathrm{A} _ {\mathrm{H}}} (x _ {\mathrm{H}}) = \mathrm{H}.\tag{18}
\]

特别地，若 H 非空，则幂等元 \(j_{H}\) 非零。

公式 (17) 与 (18) 证明：函数 \(\lambda \mapsto \mathrm{R}_{\mathrm{H}}(x,\lambda)\)（定义在 \(\mathbf{C} - \mathbf{H}\) 中）是 \(x_{\mathrm{H}}\) 相对于 \(\mathrm{A_H}\) 的预解式。

命题 16. --- 我们沿用前面的记号。设 \((\mathrm{H}_{i})_{1\leqslant i\leqslant n}\) 是 \(\mathrm{Sp}_{\mathrm{A}}(x)\) 的由 \(\Pi\) 中元素组成的一个分划。

\begin{enumerate}
\def\labelenumi{\alph{enumi})}
\item
  代数 B 典范地等同于代数 \(\mathrm{B_{H_1}}\times \dots \times \mathrm{B_{H_n}}\)
\item
  我们有 \(x_{\mathrm{H}_i}x_{\mathrm{H}_j} = 0\)（对 \(i\neq j\)），并且
\end{enumerate}

\[
x = x _ {\mathrm{H} _ {1}} + x _ {\mathrm{H} _ {2}} + \dots + x _ {\mathrm{H} _ {n}};
\]

\begin{enumerate}
\def\labelenumi{\alph{enumi})}
\setcounter{enumi}{2}
\tightlist
\item
  我们有
\end{enumerate}

\[
\mathrm{R} (x, \lambda) = \mathrm{R} _ {\mathrm{H} _ {1}} (x, \lambda) + \dots + \mathrm{R} _ {\mathrm{H} _ {n}} (x, \lambda)\tag{19}
\]

对所有 \(\lambda\in\mathbf{C}-\mathrm{Sp}_{\mathrm{A}}(x)\) 成立。特别地，若 \(\mathrm{H}\in\Pi\)，则预解式 \(\lambda\mapsto\mathrm{R}(x,\lambda)\) 在 \(\mathrm{H}\) 的邻域中等于 \(\mathrm{R}_{\mathrm{H}}(x,\lambda)\) 与一个全纯函数之和。

关系式 \(1 = j_{H_{1}} + \cdots + j_{H_{n}}\) 是 1 分解为 B 中两两正交的幂等元的分解，故代数 B 典范地等同于乘积代数 \(B_{H_{1}} \times \cdots \times B_{H_{n}}\)（A, I, p.~105, 命题 10）。

断言 \(b\)) 由函数 \(g_{\mathrm{H}_i}\) 的相应关系式得出；断言 \(c\)) 是 \(a\)) 与等式 \(\mathrm{R}(x_{\mathrm{H}},\lambda) = \mathrm{R}_{\mathrm{H}}(x,\lambda)\) 的推论。

命题 17. --- 设 \(\mu\) 是 \(\mathrm{Sp}_{\mathrm{A}}(x)\) 的一个孤立点。那么

\begin{enumerate}
\def\labelenumi{\alph{enumi})}
\tightlist
\item
  对每个 \(\lambda \in \mathbf{C} - \mathrm{Sp}_{\mathrm{A}}(x)\)，我们有
\end{enumerate}

\[
\mathrm{R} (x, \lambda) = \mathrm{R} _ {\{\mu \}} (x, \lambda) + \mathrm{R} _ {\mathrm{Sp} _ {\mathrm{A}} (x) - \{\mu \}} (x, \lambda);
\]

\begin{enumerate}
\def\labelenumi{\alph{enumi})}
\setcounter{enumi}{1}
\item
  函数 \(\lambda \mapsto \mathrm{R}_{\mathrm{Sp}_{\mathrm{A}}(x)-\{\mu\}}(x,\lambda)\) 在 \(\mathbf{C}-\mathrm{Sp}_{\mathrm{A}}(x)\) 中以及 \(\mu\) 的一个邻域中全纯；此外，函数 \(\lambda \mapsto \mathrm{R}_{\{\mu\}}(x,\lambda)\) 在 \(\mathbf{C}-\{\mu\}\) 中全纯；
\item
  我们有
\end{enumerate}

\[
\lim _ {n \to + \infty} \| (x - \mu) ^ {n} j _ {\{\mu \}} \| ^ {1 / n} = 0
\]

并且，对 \(\lambda \in \mathbf{C} - \{\mu\}\)，公式

\[
\mathrm{R} _ {\{\mu \}} (x, \lambda) = \sum_ {n = 0} ^ {\infty} (\lambda - \mu) ^ {- n - 1} (x - \mu) ^ {n} j _ {\{\mu \}}.\tag{20}
\]

前文已蕴涵断言 \(a\)) 与 \(b\))。我们来证明 \(c\))。通过把 \(x\) 换成 \(x-\mu\)，可化为 \(\mu=0\) 的情形。令 \(H = \{0\}\)；这是 \(\mathrm{Sp}_{\mathrm{A}}(x)\) 的既开又闭的子集。由公式 (18)，\(x_{H}\) 在 \(\mathrm{A}_{H}\) 中的谱是 \{0\}，因而 \(x_{H}\) 是拟幂零的，即 \(\|x^{n}j_{H}\|^{1/n}=\|(xj_{H})^{n}\|^{1/n}\) 当 \(n\) 趋于 + \(\infty\) 时趋于 0。此外，对 \(\lambda\neq0\)，在 A 中有

\[
\left(\lambda j _ {\mathrm{H}} - x _ {\mathrm{H}}\right) ^ {- 1} = \sum_ {n = 0} ^ {\infty} \lambda^ {- n - 1} x _ {\mathrm{H}} ^ {n}
\]

（I, p.~24 的定理 1, d)），由此得 (20)。

推论 1. --- 设 \(\mu\) 是 \(\mathrm{Sp}_{\mathrm{A}}(x)\) 的孤立点，\(p\) 是严格正整数。为使 \(\mu\) 成为阶为 \(p\) 的极点（对 \(x\) 的预解式而言；参见 VAR, R1, p.~30, 3.3.9），当且仅当 \((x - \mu)^{p-1} j_{\{\mu\}} \neq 0\) 且 \((x - \mu)^{p} j_{\{\mu\}} = 0\)。

推论 2. --- 设 \(\mu\) 是 \(\mathrm{Sp}_{\mathrm{A}}(x)\) 的孤立点。设 \(\Gamma\) 是一个开圆盘 \(\Delta\)（以 \(\mu\) 为中心）的定向边界，使得

\[
\operatorname{Sp} _ {\mathrm{A}} (x) \cap (\Gamma \cup \Delta) = \{\mu \}.
\]

则幂等元 \(j_{\{\mu \}}\)（与 \(x\) 和 \(\{\mu\}\) 相伴）由下式给出

\[
j _ {\{\mu \}} = \frac {1}{2 i \pi} \int_ {\Gamma} (z - x) ^ {- 1} d z.
\]

*换言之，幂等元 \(j_{\{\mu\}}\) 是在点 \(\mu\) 处 \(x.*\) 的预解式的留数。

对 \(z\in \mathbf{C} - \mathrm{Sp}_{\mathrm{A}}(x)\)，有

\[
(z - x) ^ {- 1} = \mathrm{R} (x, z) = \mathrm{R} _ {\{\mu \}} (x, z) + \mathrm{R} _ {\mathrm{H}} (x, z),
\]

其中 H = SpA(x) - \{μ\}（公式 (19)）。函数 z ↦ R\_H(x, z) 在 C-H 中以及 \{μ\} 的一个邻域中全纯（命题 17, b)），因而

\[
\frac {1}{2 i \pi} \int_ {\Gamma} \mathrm{R} _ {\mathrm{H}} (x, z) d z = 0
\]

（VAR, R2, p.~48, 11.2.5）。函数 \(z \mapsto \mathrm{R}_{\{\mu\}}(x,z)\) 是元素 \(j_{\{\mu\}}xj_{\{\mu\}}\)（属于幺代数 \(\mathrm{A}_{\{\mu\}}\)）的预解式。于是有

\[
j _ {\{\mu \}} = \frac {1}{2 i \pi} \int_ {\Gamma} \mathrm{R} _ {\{\mu \}} (x, z) d z
\]

（由应用于 \(A_{\{\mu\}}\) 以及 \(\Delta \cup \Gamma\) 的邻域中的常值函数 1 的 I, p.~75 的命题 9）。推论得证。

